\ifdefined\enmain\else
  \def\tituloconcepto{Referencias}
  % Si el capitulo se compila solo (sin main.tex), esta cabecera abre el documento.
% Antes de % Si el capitulo se compila solo (sin main.tex), esta cabecera abre el documento.
% Antes de % Si el capitulo se compila solo (sin main.tex), esta cabecera abre el documento.
% Antes de \input{cabecera} cada capitulo define \tituloconcepto.
\ifdefined\enmain\else
  \documentclass[11pt,a4paper]{article}
  \input{preambulo}
  \begin{document}
  \title{\tituloconcepto}
  \author{\autor}
  \date{\today}
  \maketitle
\fi
 cada capitulo define \tituloconcepto.
\ifdefined\enmain\else
  \documentclass[11pt,a4paper]{article}
  % Preambulo compartido por todos los documentos. Compilar con XeLaTeX/LuaLaTeX
% (o Tectonic). Si se usa Overleaf, elegir el compilador XeLaTeX.
\usepackage[spanish,es-noshorthands]{babel}
\usepackage{amsmath,amssymb}
\usepackage{booktabs,array}
\usepackage[a4paper,margin=2.5cm]{geometry}
\usepackage{xcolor}
\usepackage{listings}
\usepackage[most]{tcolorbox}
\usepackage{quantikz}
\usepackage[hidelinks]{hyperref}

% ---- Autor (editar a gusto) -------------------------------------------------
\newcommand{\autor}{Genesis Silesky}

% ---- Notacion de Dirac y atajos ---------------------------------------------
% quantikz ya define \ket y \bra; DeclareRobustCommand los redefine sin error.
\DeclareRobustCommand{\ket}[1]{\left|#1\right\rangle}
\DeclareRobustCommand{\bra}[1]{\left\langle#1\right|}
\DeclareRobustCommand{\braket}[2]{\left\langle#1\middle|#2\right\rangle}
\newcommand{\C}{\mathbb{C}}
\newcommand{\R}{\mathbb{R}}
\DeclareMathOperator{\Tr}{Tr}

% ---- Cajas de apoyo --------------------------------------------------------
% "concepto": definiciones previas que hacen falta para entender el capitulo.
\newtcolorbox{concepto}[1]{breakable,colback=blue!4,colframe=blue!55!black,
  fonttitle=\bfseries,title={Concepto previo: #1}}
% "nota": aclaraciones y advertencias.
\newtcolorbox{nota}[1]{breakable,colback=orange!6,colframe=orange!70!black,
  fonttitle=\bfseries,title={#1}}

% ---- Estilos de listados ----------------------------------------------------
\lstdefinestyle{py}{language=Python,basicstyle=\ttfamily\footnotesize,
  breaklines=true,postbreak=\mbox{\textcolor{gray}{$\hookrightarrow$}\space},
  numbers=left,numberstyle=\tiny\color{gray},frame=single,
  commentstyle=\color{green!40!black}\itshape,
  keywordstyle=\color{blue!70!black}\bfseries,
  stringstyle=\color{red!60!black},showstringspaces=false,tabsize=4,
  columns=fullflexible,keepspaces=true}
\lstdefinestyle{salida}{basicstyle=\ttfamily\footnotesize,breaklines=true,
  frame=leftline,backgroundcolor=\color{gray!8},columns=fullflexible,
  keepspaces=true}

  \begin{document}
  \title{\tituloconcepto}
  \author{\autor}
  \date{\today}
  \maketitle
\fi
 cada capitulo define \tituloconcepto.
\ifdefined\enmain\else
  \documentclass[11pt,a4paper]{article}
  % Preambulo compartido por todos los documentos. Compilar con XeLaTeX/LuaLaTeX
% (o Tectonic). Si se usa Overleaf, elegir el compilador XeLaTeX.
\usepackage[spanish,es-noshorthands]{babel}
\usepackage{amsmath,amssymb}
\usepackage{booktabs,array}
\usepackage[a4paper,margin=2.5cm]{geometry}
\usepackage{xcolor}
\usepackage{listings}
\usepackage[most]{tcolorbox}
\usepackage{quantikz}
\usepackage[hidelinks]{hyperref}

% ---- Autor (editar a gusto) -------------------------------------------------
\newcommand{\autor}{Genesis Silesky}

% ---- Notacion de Dirac y atajos ---------------------------------------------
% quantikz ya define \ket y \bra; DeclareRobustCommand los redefine sin error.
\DeclareRobustCommand{\ket}[1]{\left|#1\right\rangle}
\DeclareRobustCommand{\bra}[1]{\left\langle#1\right|}
\DeclareRobustCommand{\braket}[2]{\left\langle#1\middle|#2\right\rangle}
\newcommand{\C}{\mathbb{C}}
\newcommand{\R}{\mathbb{R}}
\DeclareMathOperator{\Tr}{Tr}

% ---- Cajas de apoyo --------------------------------------------------------
% "concepto": definiciones previas que hacen falta para entender el capitulo.
\newtcolorbox{concepto}[1]{breakable,colback=blue!4,colframe=blue!55!black,
  fonttitle=\bfseries,title={Concepto previo: #1}}
% "nota": aclaraciones y advertencias.
\newtcolorbox{nota}[1]{breakable,colback=orange!6,colframe=orange!70!black,
  fonttitle=\bfseries,title={#1}}

% ---- Estilos de listados ----------------------------------------------------
\lstdefinestyle{py}{language=Python,basicstyle=\ttfamily\footnotesize,
  breaklines=true,postbreak=\mbox{\textcolor{gray}{$\hookrightarrow$}\space},
  numbers=left,numberstyle=\tiny\color{gray},frame=single,
  commentstyle=\color{green!40!black}\itshape,
  keywordstyle=\color{blue!70!black}\bfseries,
  stringstyle=\color{red!60!black},showstringspaces=false,tabsize=4,
  columns=fullflexible,keepspaces=true}
\lstdefinestyle{salida}{basicstyle=\ttfamily\footnotesize,breaklines=true,
  frame=leftline,backgroundcolor=\color{gray!8},columns=fullflexible,
  keepspaces=true}

  \begin{document}
  \title{\tituloconcepto}
  \author{\autor}
  \date{\today}
  \maketitle
\fi

\fi
\section*{Referencias}
\addcontentsline{toc}{section}{Referencias}
\begin{itemize}
  \item M.~A.~Nielsen e I.~L.~Chuang, \emph{Quantum Computation and Quantum
    Information}, 10th Anniversary Edition, Cambridge University Press, 2010.
    (Referencia general; cap.~1 y 2 para los fundamentos.)
  \item Documentación de Qiskit y el curso \emph{Basics of quantum information}
    de IBM Quantum Learning (verificar la versión vigente de la API antes de
    usarla: cambia entre versiones).
  \item M.~Born, ``Zur Quantenmechanik der Stoßvorgänge'', \emph{Zeitschrift
    für Physik} 37, 863--867, 1926. (Regla de Born.)
  \item J.~S.~Bell, ``On the Einstein Podolsky Rosen paradox'',
    \emph{Physics} 1, 195--200, 1964.
  \item D.~M.~Greenberger, M.~A.~Horne y A.~Zeilinger, ``Going beyond Bell's
    theorem'', en \emph{Bell's Theorem, Quantum Theory and Conceptions of the
    Universe}, Kluwer, 1989. (Estados GHZ.)
  \item L.~K.~Grover, ``A fast quantum mechanical algorithm for database
    search'', \emph{Proc.\ 28th ACM STOC}, 212--219, 1996.
  \item P.~W.~Shor, ``Algorithms for quantum computation: discrete logarithms
    and factoring'', \emph{Proc.\ 35th IEEE FOCS}, 124--134, 1994.
  \item J.~Preskill, ``Quantum Computing in the NISQ era and beyond'',
    \emph{Quantum} 2, 79, 2018.
  \item Notas del curso \emph{Quantum Computing Experimental Platform}
    (materiales de la carpeta del Laboratorio de Programación Cuántica UNA),
    prácticas 1.3 \emph{Rabi Oscillations}, 1.4 \emph{Quantum Bit},
    1.6 \emph{Quantum Control} y 1.9 \emph{Quantum State Tomography}.
    (Base de los capítulos 10 a 13.)
\end{itemize}
\ifdefined\enmain\else\ifdefined\enmain\else
  \end{document}
\fi
\fi
